\pdfvariable trailerid {[<C2972026090200000000000000000000><C2972026090200000000000000000000>]}
\documentclass[10pt]{article}
\usepackage[margin=0.78in]{geometry}
\usepackage{amsmath,amssymb,amsthm,booktabs,microtype,hyperref}
\hypersetup{colorlinks=true,linkcolor=blue,urlcolor=blue}
\providecommand{\CRevisionRound}{2}
\newtheorem{theorem}{Theorem}
\newtheorem{proposition}{Proposition}
\newcommand{\C}{\mathbb C}
\newcommand{\CP}{\mathbb {CP}}
\newcommand{\dd}{\mathrm d}
\newcommand{\ii}{\mathrm i}
\title{Exact Phase, Projective, and Metric Atlases\\
for a PT-Symmetric Dimer}
\author{Route-A source-local certificate HCS-C297}
\date{2 September 2026}
\begin{document}
\maketitle

\begin{abstract}
For the autonomous balanced-gain/loss dimer
$\ii\dot\psi=H_{\gamma,\kappa}\psi$, $\kappa>0$, we give an
all-parameter dynamical classification.
\ifnum\CRevisionRound=0
The scalar identity $H^2=(\kappa^2-\gamma^2)I$ yields the exact propagator,
generic periods in the unbroken chamber, a rank-one nilpotent exceptional
point, and the attracting/repelling eigenrays of the broken chamber.
\fi
\ifnum\CRevisionRound>0
We additionally compactify the flow on $\CP^1$, separate vector and ray
periods, and prove the sharp positivity boundary for an explicit conserved
pseudo-Hermitian metric.
\fi
\ifnum\CRevisionRound>1
An independent exact receipt closes both exceptional sheets and all singular
parameter faces while enforcing a literal source/target scope firewall.
\fi
\end{abstract}

\section{Frozen system and exact propagator}
The frozen release obstruction identifier is \texttt{HEN-O281}.
Let
\[
 H=H_{\gamma,\kappa}=
 \begin{pmatrix}\ii\gamma&\kappa\\ \kappa&-\ii\gamma\end{pmatrix},
 \qquad \kappa>0,\quad\gamma\in\mathbb R,
 \qquad \ii\dot\psi=H\psi.
\]
Parity is $\mathcal P=\sigma_x$ and time reversal is componentwise complex
conjugation, so $\mathcal P H^*\mathcal P=H$.  Time is the physical
propagation parameter; it is not adjusted to external data.  Put
$\delta=\kappa^2-\gamma^2$.

\begin{theorem}[Complete three-chamber atlas]\label{thm:atlas}
The propagator is
\[
e^{-\ii tH}=\begin{cases}
\cos(\omega t)I-\ii\omega^{-1}\sin(\omega t)H,
 &\delta>0,\quad\omega=\sqrt\delta,\\[3pt]
I-\ii tH,&\delta=0,\\[3pt]
\cosh(\nu t)I-\ii\nu^{-1}\sinh(\nu t)H,
 &\delta<0,\quad\nu=\sqrt{-\delta}.
\end{cases}
\]
If $\delta>0$, a ray having nonzero components in both eigenlines has least
projective period $\pi/\omega$, while its vector representative has least
period $2\pi/\omega$.  The eigenrays themselves are stationary.  If
$\delta=0$, $H$ is nonzero rank-one nilpotent with one eigenline, and every
generalized state has linear leading growth.  If $\delta<0$, there are two
fixed rays; one attracts every generic forward ray and the other repels it,
with exponential rate $2\nu$ in their affine ratio.
\end{theorem}

\begin{proof}
Multiplication gives
\[
H^2=(\kappa^2-\gamma^2)I=\delta I.
\]
The even and odd terms of the exponential series give the three formulae.
For $\delta>0$ the eigenvalues are $\pm\omega$.  In their eigenbasis a
generic component ratio is multiplied by $e^{2\ii\omega t}$, proving the ray
period; equality of both vector phases first occurs at $2\pi/\omega$.  For
$\delta=0$, nonzero determinant-zero $H$ has rank one and the exponential
terminates.  For $\delta<0$ the eigenvalues are $\pm\ii\nu$, so the evolution
multipliers are $e^{\pm\nu t}$ and their ratio proves the final assertion.
\end{proof}

This proof is uniform at the exceptional point: it does not diagonalize there
or suppress the generalized eigenvector.  It also prevents a common factor
of two error between vector and ray periods.

\ifnum\CRevisionRound>0
\section{Projective flow and conserved metrics}
On the affine chart $\psi_1\ne0$, set $z=\psi_2/\psi_1$.  Direct
differentiation gives the Riccati equation
\begin{equation}\label{eq:riccati}
 \dot z=\ii\kappa(z^2-1)-2\gamma z.
\end{equation}
The reciprocal chart extends it across infinity to a complete holomorphic
vector field on $\CP^1$.  The complex quadratic discriminant of its fixed
point equation is $-4\delta$; the matrix characteristic discriminant is
$+4\delta$.  Thus the two fixed rays merge precisely on
$|\gamma|=\kappa$.  These opposite signs are conventionally harmless but
must not be conflated.

\begin{theorem}[Sharp metric boundary]\label{thm:metric}
With
\[
 Q=\sigma_x,\qquad
 \eta=I+\frac{\gamma}{\kappa}\sigma_y,
\]
one has $H^\dagger Q=QH$ and $H^\dagger\eta=\eta H$.  Both quadratic
forms are conserved.  The form $Q$ is indefinite for every parameter, while
$\eta$ is positive definite exactly when $|\gamma|<\kappa$, singular
positive semidefinite at $|\gamma|=\kappa$, and indefinite when
$|\gamma|>\kappa$.
\end{theorem}

\begin{proof}
The intertwining identities follow by direct Pauli-matrix multiplication.
For either matrix $M$, differentiation gives
\[
 \frac{\dd}{\dd t}(\psi^\dagger M\psi)
 =\ii\psi^\dagger(H^\dagger M-MH)\psi=0.
\]
The eigenvalues of $\eta$ are
$1+|\gamma|/\kappa$ and $1-|\gamma|/\kappa$, proving the signature claim.
\end{proof}

The standard Euclidean norm is not substituted for these forms:
\[
 \frac{\dd}{\dd t}\|\psi\|^2
 =2\gamma\,\psi^\dagger\sigma_z\psi.
\]
It is generally not conserved.  In the unbroken phase the two real-energy
eigenlines have opposite $Q$-Krein signs; their collision makes the normalized
positive metric singular at the exceptional point.
\fi

\ifnum\CRevisionRound>1
\section{Boundary closure and independent evidence}
Table~\ref{tab:boundary} records the faces excluded from abbreviated phase
diagrams.
\begin{table}[h]
\centering
\small
\begin{tabular}{@{}ll@{}}
\toprule
Face & Exact outcome\\
\midrule
$\gamma=0$ & Hermitian $H=\kappa\sigma_x$ and $\eta=I$\\
$\gamma=\pm\kappa$ & two nilpotent sheets, one eigenline on each\\
$\kappa\downarrow0$, $\gamma\ne0$ & uncoupled amplification/decay; no positive displayed $\eta$\\
$(\kappa,\gamma)=(0,0)$ & zero generator, not an exceptional point\\
eigenray initial data & projectively stationary, not generic-periodic\\
$\psi=0$ & fixed vector with no point in $\CP^1$\\
\bottomrule
\end{tabular}
\caption{Boundary atlas.}\label{tab:boundary}
\end{table}

A canonical exact receipt enumerates all 168 integer cells
$1\le\kappa\le8$, $-10\le\gamma\le10$: 64 are unbroken, 16 exceptional,
and 88 broken.  A producer-independent checker reconstructs them with 6,475
assertions.  A separate symbolic route contributes 516 identities and grid
checks; two fresh paths reproduce the same exact evidence byte.  All 52 hostile
schema, repaired-hash, duplicate-key, obstruction-ID, scope, and YAML
graph-sharing attacks
are rejected.  The finite rectangle is regression evidence; Theorems
\ref{thm:atlas}--\ref{thm:metric} are all-parameter algebraic statements.

\section{Route-A result, ownership, and limitations}
The foundational PT-symmetric setting belongs to Bender and
Boettcher~\cite{BB1998}; positive-metric pseudo-Hermiticity is treated by
Mostafazadeh~\cite{M2002}, and balanced optical gain/loss has an experimental
owner in Rueter et al.~\cite{R2010}.  Our result is a source-model closure and
does not claim those mechanisms as literature originality.

Under evaluator v0.2.0 the strict tuple is
\[
(\mathrm{A0\_FAIL},\mathrm{A1\_WEAK},\mathrm{A2\_FAIL},
 \mathrm{A3\_FAIL},\mathrm{A4\_FORMAL\_HINT}),
\]
with overall verdict \texttt{ROUTE\_A\_REJECTED}.  A1 is only a clean
projective family, not an isolated arithmetic periodic-orbit bridge.  A4 is
only a finite-dimensional, parameter-dependent metric hint.  No fixed
Hilbert-space self-adjoint realization across the exceptional point is
asserted.

The literal scope is \texttt{NO\_BAD\_EULER\_OR\_ROOT\_NUMBER}.  We construct
no arithmetic local data, Euler factors, root numbers, automorphy, target
divisor or functional equation, target zero match, Hilbert--P\'olya operator,
or Route-B input.
\fi

\begin{thebibliography}{9}
\bibitem{BB1998} C.~M. Bender and S.~Boettcher,
``Real spectra in non-Hermitian Hamiltonians having PT symmetry,''
\emph{Phys. Rev. Lett.} 80 (1998), 5243--5246.
\bibitem{M2002} A.~Mostafazadeh,
``Pseudo-Hermiticity versus PT symmetry,''
\emph{J. Math. Phys.} 43 (2002), 205--214.
\bibitem{R2010} C.~E. Rueter et al.,
``Observation of parity-time symmetry in optics,''
\emph{Nature Phys.} 6 (2010), 192--195.
\end{thebibliography}
\end{document}
